\subsection{All terminals beginning at three, at both remaining roots}
\label{sec:odd-three:roots}

\begin{theorem}\label{thm:odd-three:roots}
Let $3\le a<b$ or $3\le a<b<c$. The internal root is, respectively,
$pW_a\bar f_b$ or $f_a h_bf_c$. For a terminal residual $(3,t)$,
$t\ge4$, use the actual direction
\[
 F_{t,L}=\begin{cases}
 h_3f_t,&L=1\text{ or }t=4,\\
 \bar f_3W_t,&t\ge5,\ L\ge2.
 \end{cases}
\]
For every finite $t$, every $L\ge1$, and the whole host domain,
the resulting complete message $H_0$ has
\begin{equation}\label{eq:odd-three:fourth}
 \E H_0^4\le e^{-6d}.
\end{equation}
On $0<d\le1/64$, its proof supplies a pointwise certificate
\begin{equation}\label{eq:odd-three:vector}
 \log H_0\le a_0+BV_0+W_0,\qquad
 \|V_0\|<\tfrac72\sqrt d,
\end{equation}
with $W_0=0$ or an actual already paid $p$-curvature term.
All later ordered root pairs, genuine equal pairs and every independent
outer length are therefore covered by
Lemmas~\ref{lem:odd-common:extra-pairs} and
\eqref{eq:odd-common:outlet}.
\end{theorem}
\begin{proof}
The separate terminal directions are fixed by the original $t,L$.
All intermediate means, constants and vectors below use the original
measure. We give one actual terminal packet and one four-row payment
for both root classes.

\paragraph{The packet for $t\ge5$.}
Write $G_{t-2}=A+zC$, $w=1-A\in[0,x]$, and
$b_0=\beta_{t-1}=\int C\,d\nu$. The actual two directions have
\[
 F_{t,1}\le e^{R_+-\Gamma},\quad
 F_{t,L}\le e^{R_--\Gamma}\ (L\ge2),\qquad
 R_\pm=(b_0-xw\mp zxC)u,
 \quad\Gamma=d+\beta_t+b_0\ge3d-2\delta_2.
\]
These are the arm-sum identities of part~\ref{part:C160.4}.
Their first positive rows use the same actual remainders.
For $R_-$, the complete finite calculation in Section~\ref{proof:Q5-C160}
gives the following data on the entire domain:
\[
 \begin{gathered}
 \|R_-\|<\tfrac12\sqrt d,\quad
 -\delta_2\le\delta<d/24,\quad
 -3d/4\le\epsilon=\delta+D_1\le0,\\
 (D_1)_+\le\delta_3,\qquad
 \beta_j-D_{j-1}\ge9d/16\ (j\ge3).
 \end{gathered}
\]
Here $D_i=\E B^iR_-$ and the Taylor coefficient is $5/4$.
For $R_+$ at $L=1$, use coefficient $3/2$ and its actual unit
certificate. The finite propagation is the bare exponent packet
\eqref{eq:oc:fourth-duhamel}.

Put $n=L-1$ in the long row and $h=9/16$.
In the full-constant convention its propagated coordinates are
\begin{equation}\label{eq:odd-three:actual-rows}
\begin{array}{c|c|c}
 L&Q&\gamma\one\text{ retained in the full vector}\\\hline
 1&r_+&0\\
 2&1+zr_-&\delta\one\\
 3&G_2+\delta+z^2r_-&\epsilon\one\\
 \ge4&G_n+z^nr_-&-hd\one.
\end{array}
\end{equation}
The mean and scalar in each row are obtained from this actual upper,
including the same $D_i$; no alternate full polynomial is substituted.
At the auxiliary root $P_0$, the complete baseline losses are
\begin{equation}\label{eq:odd-three:baselines}
 25/4,\qquad8,\qquad13,\qquad3385/384.
\end{equation}
The first two are the actual unit and two-step calculations of
Section~\ref{proof:Q5-C160}. For the long row its complete calculation
has loss $1043/128$ after using the weaker root loss $5d/6$.
At this same $v=3$ root the loss is exactly $d$ and its constant
$\beta_2$ is zero, saving $2d/3$ without changing the full vector.

For the third row a direct specialization is substantially shorter.
The nonconstant old vector $(Q-B)u$ has real coordinate
$1+\delta-b_0x+x^2w$ and imaginary coefficient $-x^2C$.
The real part is positive and at most
$1+1/72+\theta^3<13/12$, so its squared norm is below $6d/5$.
Its actual mean $\int x^2C\,d\nu$ is nonnegative.
Since $\epsilon\le0$, the cross with $\epsilon\one$ helps; its
square costs at most $3d/16$. The complete norm is thus below
$7d/5$. The scalar plus mean is at most
$[-6+2\theta+\theta^2+2\theta/3]d$. Consequently its loss exceeds
\[
 24-8\theta-4\theta^2-8\theta/3-119/20
       =50137/3840>13.
\]

\paragraph{The two root transfers for this same packet.}
For the minimum-one root, use
Lemma~\ref{lem:odd-complete:minimum-one} in the full-constant
convention. At $L=1$ the signed imaginary part gives a stronger
direct comparison: its nonconstant norm increment is at most
\[
 \int(-x+x^2+3x^3)\,d\nu\le0.
\]
Indeed the real increment is at most $x^2+2x^3$, while the imaginary
increment is at most $-x+x^3$. The scalar gain is $25d/8$ and the
new curvature costs at most $10d/9$, giving a further improvement
$145d/72$. For $a\ge4$, $\Psi(1,1/2)<0$ suffices.

At $L=2$, $A_Q\ge0$, $D_Q\le1$, $\|Qu\|\le(6/5)\sqrt d$,
and $\delta\ge-\theta d$. The worst full-constant expense is
\[
 \Phi(1,0,1/3)+(17/6)\theta
       =31177/36864<17/20;
\]
the $a\ge4$ expense is negative.
At $L=3$, take $\bar D=1$, $R=7/5$, $\epsilon\ge-3d/4$;
the two expenses are below $2d$ and $d$, respectively.
For the long row, $\bar D=5/4$, $R=6/5$, $\gamma=-9d/16$;
the two expenses are below $(11/5)d$ and $d$.
These bounds keep the complete constant vector in every row.

For the positive root use
Lemma~\ref{lem:odd-complete:positive}, again in its full-constant
convention. At $L=1$ the only adverse overshoot is the negative
imaginary part of size at most $x$, costing $17d/72<d$.
At $L=2$ its size is at most $x^2$, costing $17d/168$.
The signed constant correction is at most $(17/8)d^2$, since
$\delta\ge-\delta_2$ and $\delta_+<d/24$.
Together with the saved $d$, this row improves.

At $L=3$, the real overshoot is at most $d/24+x^3$ and the
imaginary coordinate is in $[0,1]$. Its charge is below
$(119/288)d^2+(119/168)d<17d/20$.
The complete $\epsilon\one$ correction is below
$(15/4)d^2\le5d/4$. The increment is therefore below $11d/10$.
For the long row $n=3$ has coordinates
$1-x+x^3C$, $1-b_0x+x^2w$, with only imaginary overshoot $x^3$.
Its charge is below $17d/112<d/5$.
For every $n\ge4$ the coordinate perturbations are at most
$x^2(1/3+x)$; they lie between the adjacent finite alternating
prefixes, because $(4/3)x+x^2<1$. Thus they stay in $[0,1]$.
The negative constant $-hd\one$ costs less than $3d^2\le d$,
paid by the root saving. The long increment is below $d/5$.

The resulting simultaneous payment table is
\[
\begin{array}{c|c|c|c}
 L&\text{actual old full loss}&\text{minimum-one expense}/d
       &\text{positive-root expense}/d\\\hline
 1&25/4&0&0\\
 2&8&17/20&0\\
 3&13&2&11/10\\
 \ge4&3385/384&11/5&1/5.
\end{array}
\]
Every loss minus its indicated expense is greater than six.

\paragraph{The adjacent terminal $t=4$.}
Keep the actual $F=h_3f_4$ for every $L$. Its one packet is
\[
 \begin{gathered}
 \Gamma=3d,\quad R=(d-zx)u,\quad
 \delta=\delta_2+(5/4)(d^3+\delta_4),\\
 \epsilon=\delta-d^2\ge0,\quad
 \delta<101d/160,\quad\delta^2<2d^2/5.
 \end{gathered}
\]
Every finite loss is at least $d/2$, by the signed moments of
Section~\ref{proof:Q5-C170}. Use the same four rows as
\eqref{eq:odd-three:actual-rows}, with $r=d-zx$, $h=1/2$,
but with the actual curved convention.
The auxiliary complete baseline losses can be taken
\[
 301/40,\qquad9,\qquad14,\qquad547/60.
\]
Here the unit and long rows are the original complete calculations.
For $L=2$, the actual $v=3$ scalar plus mean is at most
$-699d/160$; its nonconstant squared norm is below $16d/9$,
and its curvature costs at most $4d/9$. The loss is thus greater
than $699/40-68/9-4/9=379/40>9$.
At $L=3$ the nonconstant coordinates are
$1-dx+\delta$, $x^2$, with squared norm below $3d/2$.
The scalar plus mean is at most $-859d/160$ and the actual
$\epsilon$ curvature costs at most $4d/9$. This gives loss
$859/40-51/8-4/9=1319/90>14$.

For the minimum-one root, the four curved expenses are at most
$0,0,3/5,21/10$, using respectively
$(\bar D,R)=(1,6/5),(1,6/5),(5/4,3/2),(3/2,8/5)$.
All resulting losses exceed six on the full domain.
For the positive root first take $d\le1/8$.
The first row improves as above. In the second row, real overshoot
$x^2$ costs at most $119d/72$, and its actual curved correction
is less than $d$. After the saved $d$, its increment is less than
$2d$. In the third row the two overshoots cost less than $9d/8$;
its new virtual coordinate is nonnegative, so $\epsilon\E V_*\le0$.
Only $11\epsilon^2/12\le11d/240$ remains from the curved correction.
The increment is below $2d$.
The long tail has norm below $(8/75)\sqrt d$; the integral
comparison costs $47d/225$. Its curved correction is below
$9d^2\le9d/8$, again leaving an increment below $2d$.
The baseline losses therefore give at least six in all four windows.

On $d\ge1/8$, the actual arm sum is
$h_3+f_4=2-3d+du+B^3u$.
The centered-row inequality gives
$|du+B^3u|\le\sqrt{d^2+\delta_3}\le2d$, since
$\delta_3\le\theta^2d\le8\theta^2d^2<3d^2$.
Thus this same actual $F=h_3f_4$ satisfies $F\le e^{-d}$.
Thus $T^LF\le e^{-d}\bar f_L$, and
Corollary~\ref{cor:odd-complete:bare} gives positive-root loss at
least $4+387/128>6$ with the same directions.

\paragraph{Actual full vectors and later arms.}
On $d\le1/64$, the packet rows above have
$\|Q(B)u\|\le(3/2)\sqrt d$ and
$|\gamma|/\sqrt d\le3/32$.
The minimum-one root adds a nonconstant precursor of norm at most
$\sqrt d/2$; its constant or curve is the one used in the proof.
The positive root adds its actual anchor precursor of norm below
$(6/5)\sqrt d$ and its pair precursor below $(3/8)\sqrt d$.
Consequently the complete vector is below
$(3/2+6/5+3/8+3/32)\sqrt d<(7/2)\sqrt d$.
This proves \eqref{eq:odd-three:vector} and permits the existing
paid restoration of every later ordered root pair and independent
outer connector. Equal pairs are restored by positivity.
At $d=0$ leaf integration gives one, with no requirement that $B=0$.
\end{proof}
